The pilot ran on a preliminary frame built from a leading portion of the bulk
opinion table, before the full pass was committed to. Its purpose was to find out
whether the guideline could be applied at all, whether the controls behaved, and
which strata were worth sampling heavily. Sixty items were annotated. Pilot items
sit outside the benchmark, and every pilot label is excluded from the results.

\subsection{Pilot results} Labels fell out by stratum as follows.
\textsc{ctrl\_hold} 12 \textsc{decided} and 0 \textsc{unresolved};
\textsc{plain} 19 \textsc{unresolved} and 1 \textsc{decided};
H2 8 \textsc{unresolved} and 2 \textsc{decided};
H1 5 \textsc{unresolved}, 5 \textsc{decided}, and 5 \textsc{unclear}.

Three things follow. The holding-cue control behaves as intended and stays clean.
The \textsc{plain} stratum is nearly degenerate, since a trigger with no
structural complication almost always does mean the court left the issue open,
which is why a keyword rule scores as well as it does and why that score carries
little information. The H1 stratum is close to an even split, so it is where the
task is genuinely difficult, and the sample was reallocated toward it. This
reallocation changes only how many items are drawn per stratum. Label definitions
and stratum definitions were fixed before the pilot and held throughout.

\subsection{Protocol revision} Version 1.0 of the guideline lacked a rule for a
passage whose only trigger sits inside a parenthetical or a description of
another court's work, on a question this opinion itself leaves aside. A citation noting that the Supreme Court declined to decide whether an
airport user fee was excessive is silent on the disposition of the Federal
Circuit opinion that cites it, and forcing a choice between
\textsc{decided} and \textsc{unresolved} would have manufactured a label. Rule
L1-9 routes these to \textsc{unclear}, using an existing label rather than adding
one. This is the single revision the design permits, it was made once, before any
main-sample annotation, and the changelog is reproduced in
Appendix~\ref{app:guide}.

\subsection{Pipeline defects} These concern candidate construction rather than
the guideline. The pattern flagging a trigger as plausibly describing another
institution was case-sensitive, matching ``the district court'' while missing
``the District Court,'' which appellate opinions write about as often. It now
also recognizes ``the court held'' and its near neighbours. The annotator was
shown a later holding sentence only for items already flagged H2, though Rule
L1-3 turns on precisely that passage, so an H1 item whose opinion later resolved
the issue was unanswerable rather than hard. The sentence is now shown whenever
one exists outside the displayed window. And ``our holding'' was a control
locator, but in the pilot it pointed at a named prior case rather than the
present disposition every time it fired, so it was dropped.

\subsection{Scope of the pilot} The pilot was used to allocate the sample and to
fix the three defects above, and it left every label boundary untouched. The
trigger dictionary was frozen before the pilot and is unchanged. In particular,
the pilot surfaced two recurring dictionary false positives and left both in
place. One is the statutory-scope reading of \emph{reach}, as in ``the shipment
clause does not reach shipments between the continental States,'' which is a
holding about a statute rather than a non-resolution. The other is a line-break
artefact in the source text that splits \emph{preserving} into \texttt{P} and
\texttt{reserving}, which then matches T11. Both stay in, are labelled
\textsc{decided} by the guideline, and are counted against the dictionary in
\S\ref{sec:models}.
